%% ==========================================================================
%%  FINAL REPORT --- ABSTRACT
%%
%%  House style for every final_* file and every file final.tex inputs: no em
%%  dashes and no semicolons in prose. Quotations, table cells, TikZ and
%%  citation ranges are exempt. Written with the Write/Edit tools only.
%% ==========================================================================
\begin{abstract}
\noindent
A recommender learns from what its users consume, so a private recommender has
to protect two things: the ratings the model is trained on, and the items a
user fetches once the recommendations arrive. \nudgeplain{}, the strongest
recent system for the first, trains a matrix factorisation across three servers
under replicated secret sharing, publishes the item embeddings in the clear, and
leaves the fetch to ``other means''. \pirsonaplain{} fetches privately but
trains with four parties over a Boolean circuit.

We build \sys{}, the composition of the two: \nudgeplain{}'s power-iteration
training under 2-of-3 replicated sharing, followed by two-server private
information retrieval from a distributed point function, with one implementation
of that function serving both the training protocol's comparison gate and the
read. On MovieLens-100K, composing costs exactly the sum of the two halves, to
the byte. The spec-faithful normaliser, which reveals nothing, costs between
0.0018 and 0.0041 of nDCG@20 against one that reveals a norm, and 5.3 times the
rounds. Truncation, not the comparison gate, is where the rounds go. Without the
private read, ten observed fetches recover a user's private embedding to a cosine
of 0.84, and neither decoy traffic nor \nudgeplain{}'s own differential privacy
brings the attack down to what popularity alone achieves.

We also report where the composition is weaker than its parts. On a wide-area
link, downloading the whole catalogue beats our unbatched private read by a
factor of 4.6. Until we added a one-round check, a malicious client could cast a
vote of arbitrary weight in the next model, and sixteen such clients destroyed
it. The check stops inflated weight but not a forged key that redistributes it.
Security is semi-honest with at most one corrupted server, and it is argued by a
simulation sketch, not proven.
\end{abstract}
